% BHU PYQ -- Transcribed & Corrected
% Paper: CHEMN41 : Physical Chemistry (Mid-Semester)
% Exam: B.Sc. (Semester IV) Mid-Semester Examination, 2025-26 (NEP)
% Department: Chemistry

\documentclass[11pt,a4paper]{article}
\usepackage[a4paper,top=2.5cm,bottom=2.5cm,left=2.8cm,right=2.8cm,headheight=14pt]{geometry}
\usepackage{amsmath,amssymb}
\usepackage[T1]{fontenc}
\usepackage[utf8]{inputenc}
\usepackage{lmodern,microtype,fancyhdr,enumitem,hyperref,lastpage,parskip}

\hypersetup{
    colorlinks=true,
    urlcolor=black,
    linkcolor=black,
    pdftitle={CHEMN41: Physical Chemistry (Mid-Sem) -- BHU Sem IV 2025-26 (NEP)}
}

\pagestyle{fancy}
\fancyhf{}
\renewcommand{\headrulewidth}{0.4pt}
\renewcommand{\footrulewidth}{0.4pt}
\fancyhead[L]{\small   \textit{Banaras Hindu University}}
\fancyhead[R]{\small   \textit{CHEMN41: Physical Chemistry (Mid-Sem 2025-26)}}
\fancyfoot[C]{\small Page  \thepage\ of \pageref{LastPage}}


\newlist{parts}{enumerate}{2}
\setlist[parts,1]{label=(\alph*),leftmargin=*,itemsep=6pt,topsep=2pt}
\setlist[parts,2]{label=(\roman*),leftmargin=*,itemsep=4pt,topsep=2pt}

\newcommand{\pts}[1]{\hfill   \textit{[#1]}}


\begin{document}


\begin{center}
  {\large\bfseries BANARAS HINDU UNIVERSITY}\\[2pt]
  {\large\bfseries DEPARTMENT OF CHEMISTRY}\\[2pt]
  {\normalsize B.Sc. --- Semester IV (NEP) Mid Semester Examination, 2025-26}\\[6pt]
  \rule{0.8\linewidth}{0.5pt}\\[6pt]
  {\large\bfseries PHYSICAL CHEMISTRY}\\[4pt]
  {\large\bfseries Paper Code: CHEMN41 (Mid-Semester)}\\[4pt]
  {\normalsize Time: 1 Hour\hfill Maximum Marks: 20}
\end{center}

\rule{\linewidth}{0.4pt}

\smallskip



\noindent   \textit{Instructions: Attempt all questions. The figures in the right-hand margin indicate marks.}

\bigskip




\noindent   \textbf{Question 1.}

\begin{parts}
  \item Write the expression of chemical potential of a gas in real gas mixture, solvent and solute in non-ideal liquid solution.\pts{2}
  \item State physical significance of chemical potential.\pts{1}
  \item State the relation between fugacity of gas with its pressure at (i) Boyle temperature (ii) very high pressure (iii) very low pressure.\pts{2}
  \item Derive the Gibbs-Duhem equation and its significance.\pts{3}
  \item State Raoult's and Henry's law. On what type of solution, each of the law is applicable.\pts{2}
  \item Why colligative properties of a solution are only dependent on the number of solute particles present not on their chemical identity?\pts{1}
  \item Using a chemical potential vs temperature plot, show that the boiling point of a liquid increases in solution compared to that of the pure liquid.\pts{2}
  \item Define Curie and Becquerel units of radioactivity.\pts{2}
  \item The activity of radioactive isotope reduced by 25 percent after 100 minutes. Calculate the decay constant and half-life of the isotope.\pts{2}
  \item Using the Bethe notation for the nuclear reactions, complete the following reactions:\pts{2}
  
\begin{parts}
    \item ${}^{197}\mathrm{Au}\,(n, \gamma) \longrightarrow$
    \item ${}^{27}\mathrm{Al}\,(\gamma, n) \longrightarrow$
  \end{parts}
\end{parts}

\vfill
\rule{\linewidth}{0.4pt}

\begin{center}\small
\href{https://bhu-exam.vercel.app/}{Click here to visit our website}\\[4pt]
\href{https://me-aryan.vercel.app/}{Click here to visit our contributor}
\end{center}

\end{document}
